\documentclass[11pt,letterpaper]{article}
\usepackage[T1]{fontenc}
\usepackage{lmodern,amsmath,amssymb,amsthm,microtype}
\usepackage[margin=1in]{geometry}
\usepackage[hidelinks]{hyperref}
\newtheorem{theorem}{Theorem}
\title{Excluding one and two isolated points in the height-six Frankl model}
\author{Aafi Khan Malek\\ \small Independent researcher, Ahmedabad, India\\
\small\href{mailto:aafimalek2032@gmail.com}{aafimalek2032@gmail.com}}
\date{28 November 2026}
\begin{document}
\maketitle

Let \(\mathcal F\) be a finite nonempty union-closed family of nonempty
sets of height at most six: it contains no strict inclusion chain of seven
members. Call it weak if each ground point occurs in at most half
of its members, and suppose it has minimum ground cardinality among
weak families. Write \(p\to q\) for the proper forcing relation:
\(p\ne q\) and every member containing \(p\) contains \(q\). A point is
isolated when it is incident to no proper forcing edge.

\begin{theorem}
The forcing graph of a ground-minimal weak family of height at most six
has at least three isolated points.
\end{theorem}

\paragraph{Previously verified structural facts.}
In the presence of forcing the directed graph consists of outward stars
and isolated points. Let \(S\) be the set of its star centers and let
\(k=|S|\). The partition gives at least \(2k+t\) ground points,
where \(t\) is the number of isolates. The
\href{https://github.com/Aafimalek/novel-discoveries/tree/043d266655932e3badfe48220085216e90065dd6/frankl-height-six}
{preceding verification} gives a pair design \(\mathcal B\) on \(S\):
each block has at least two points and each pair of centers occurs
in exactly one block. For distinct centers \(p,q\) in a block \(B\),
write \(F_{p,q}\) for the two-leaf avoidance fiber and \(W(p,q)\)
for the number of blocks avoiding both centers. The checked
count and height-three signature inequalities give
\[
  3k+t-1+W(p,q)\le |F_{p,q}|\le H(c+t),
  \qquad H(a)=1+a+\binom a2,
\]
where \(c\) is the number of distinct occurrence signatures of
centers outside \(B\) in the compressed fiber. For any \(v\in B\),
blocks through \(v\) other than \(B\) partition the outside centers into
groups: take their members other than \(v\). Each group has one
occurrence signature by the checked block-link lemma. Thus
\(c\le g_v-1\), where \(g_v\) is the number of blocks through \(v\).

\paragraph{Minimum block degree and number of blocks.}
Set \(g=\min_{v\in S}g_v\) and \(N=|\mathcal B|\), with \(k\ge2\).
Each center's partners are partitioned by the blocks through it,
so \(1\le g\le g_v\le k-1\). Let \(Q=\sum_{A\in\mathcal B}|A|\).
Counting incidences and then ordered pairs gives
\[
 Q\ge kg,\qquad
 \sum_{A\in\mathcal B}|A|^2=k(k-1)+Q.
\]
Cauchy--Schwarz and monotonicity of \(x^2/(k(k-1)+x)\) for
\(x\ge0\) imply
\[
 N\ge\frac{Q^2}{k(k-1)+Q}\ge\frac{kg^2}{k-1+g}
   \ge\frac{g(g+1)}2.
\]
The last inequality is equivalent to \((g-1)(k-g-1)\ge0\)
after canceling the positive factor \(g\).

\paragraph{Contradiction.}
Choose a center attaining \(g\) and any block \(B\) containing it.
Pick two distinct centers \(p,q\) in \(B\). By uniqueness of their
common block,
\(W(p,q)=N-g_p-g_q+1\ge N-2(k-1)+1\). Hence
\[
 |F_{p,q}|\ge N+k+t+2.
\]
When \(1\le t\le2\), the block through the minimum-degree center
gives \(c+t\le g+1\), and therefore
\[
 |F_{p,q}|\le H(g+1)=\frac{g(g+1)}2+g+2\le N+k+1,
\]
which excludes one and two isolates when there are at least two centers.

If there is forcing and \(k=1\), the previously checked
unique-source theorem requires at least three isolated points.
If there is no forcing, every ground point is isolated; a ground
of size one or two cannot support a weak family. For size one,
its point occurs in every member. For size two, the two-element
ground set belongs to the family by union closure, and with at
most three nonempty subsets available, at least one point has
strict majority (indeed if both have frequency at most half, the
sum of the two frequencies cannot exceed the family size, which
is contradicted by the two-element member). The previously verified
projection theorem \texttt{isolated\_nonempty\_of\_weakMinimalGround}
excludes zero isolates, so the lower bound is three.

\paragraph{Scope and contributions.}
This argument is analytic and uses no bounded search. The structural
lemmas invoked above have Lean proofs; the new composition is
given here in ordinary mathematical form. Families with three or
more isolated points are not excluded by this theorem. Aafi Khan
Malek selected and directed the research. ChatGPT assisted with
proof development, formalization and writing. Codex was not used.
This note has not been peer reviewed.
\end{document}
